\pdfoutput=1
\documentclass[a4paper,11pt]{amsart}
\usepackage[hmarginratio={1:1},vmarginratio={1:1},lmargin=60.0pt,tmargin=60.0pt]{geometry}

\allowdisplaybreaks

\usepackage[numbers]{natbib}
\usepackage{latexsym,mathtools,amssymb,stmaryrd}
\usepackage{enumitem}
\usepackage[dvipsnames]{xcolor}
\usepackage{diagbox,makecell,caption}
\usepackage{tikz}
\usepackage{pgfplots}
\usepgfplotslibrary{polar}
\pgfplotsset{compat=newest}
\usepackage{multirow}
\usetikzlibrary{cd,decorations.markings,arrows.meta,positioning,calc,tqft}
\usepackage{aliascnt,etoolbox}
\usepackage[hypertexnames=false]{hyperref}
\usepackage{bookmark}
\usepackage{listings}

\captionsetup[table]{skip=10pt}

\setlist[enumerate]{itemsep=0.15cm,label=\emph{\upshape(\alph*)}}
\setlist[enumerate,2]{itemsep=0.15cm,label=\emph{\upshape(\roman*)}}

% colors
\definecolor{spinach}{RGB}{46,139,87}
\definecolor{tomato}{RGB}{255,99,71}
\definecolor{orchid}{RGB}{143,40,194}
\definecolor{lava}{RGB}{207,16,32}
\definecolor{cream}{RGB}{255,253,208}
\definecolor{verdigris}{RGB}{67,179,174}
\definecolor{mydarkblue}{RGB}{10,10,170}

% general notation
\let\emph\relax
\DeclareTextFontCommand{\emph}{\bfseries\em}
\renewcommand{\dots}{\text{...}}

\newcommand{\eg}{\text{e.g.}}
\newcommand{\ie}{\text{i.e.}}
\newcommand{\Z}{\mathbb{Z}}
\newcommand{\K}{\mathbb{K}}
\newcommand{\C}{\mathbb{C}}
\newcommand{\End}{\mathrm{End}}
\newcommand{\Hom}{\mathrm{Hom}}
\newcommand{\id}{id}

\newcommand{\paraa}{\alpha}
\newcommand{\parab}{\beta}
\newcommand{\parac}{\gamma}

\newcommand{\cob}{\mathbf{2Cob}}
\newcommand{\mcob}[1]{\mathbf{MCob}_{#1}}
\newcommand{\mucob}[3]{\mathbf{MCob}_{#1, #2, #3}}

% tikz
\tikzset{
anchorbase/.style={baseline={([yshift=#1]current bounding box.center)}},
anchorbase/.default={-0.5ex},
tinynodes/.style={font=\tiny,text height=0.25ex,text depth=0.05ex},
mor/.style={line width=0.75,color=black,fill=cream},
usual/.style={line width=1.2,color=black},
match/.style={line width=1.2, densely dotted, color=verdigris},
ribbon/.style={double, double distance=4pt, line width=0.8pt, black},
ribbon boundary/.style={line width=0.8pt, black},
vertex/.style={circle, draw, fill=white, minimum size=14pt, line width=0.8pt},
hol/.style={
decoration={markings,
post length=0.25mm,
pre length=0.25mm,
mark=at position #1 with {\node[circle,radius=0.15cm,inner sep=-1.2pt,draw,color=black,fill=white]{};}
},
postaction={decorate}
},
mob/.style={
decoration={markings,
post length=0.25mm,
pre length=0.25mm,
mark=at position #1 with {\node[circle,radius=0.15cm,inner sep=-1.2pt,draw, color=tomato,fill=tomato]{};}
},
postaction={decorate}
},
dot/.style={
decoration={markings,
post length=0.25mm,
pre length=0.25mm,
mark=at position #1 with {\node[circle,radius=0.5cm,inner sep=-1.5pt,color=black,fill=black]{};}
},
postaction={decorate}
},
dot/.default=1
}

\let\oldlightning\lightning
\renewcommand{\lightning}{\textcolor{tomato}{\pmb{\oldlightning}}}

% theorems
\usepackage{amsthm}
\newtheorem*{Theorem}{Theorem}
\newtheorem*{lemma}{Lemma}

\setcounter{tocdepth}{1}
\setcounter{secnumdepth}{3}
\hypersetup{
pdftoolbar=true,
pdfmenubar=true,
pdffitwindow=false,
pdfstartview={FitH},
pdftitle={Additional material},
pdfauthor={Daniel W. Collison and Daniel Tubbenhauer},
pdfsubject={},
pdfcreator={Daniel W. Collison and Daniel Tubbenhauer},
pdfproducer={Daniel W. Collison and Daniel Tubbenhauer},
pdfkeywords={},
pdfnewwindow=true,
colorlinks=true,
linkcolor=mydarkblue,
citecolor=teal,
filecolor=magenta,
urlcolor=orchid,
linkbordercolor=lava,
citebordercolor=teal,
urlbordercolor=orchid,
linktocpage=true
}

\newcommand{\nnfootnote}[1]{%
\begin{NoHyper}
\renewcommand\thefootnote{}\footnote{#1}%
\addtocounter{footnote}{-1}%
\end{NoHyper}
}

% code listings
\definecolor{codegreen}{rgb}{0,0.6,0}
\definecolor{codepurple}{rgb}{0.58,0,0.82}
\definecolor{backcolor}{rgb}{0.95,0.95,0.92}

\lstdefinestyle{mystyle}{
backgroundcolor=\color{backcolor},
commentstyle=\color{codegreen},
keywordstyle=\color{magenta},
stringstyle=\color{codepurple},
basicstyle=\ttfamily\footnotesize,
breakatwhitespace=false,
breaklines=true,
captionpos=b,
keepspaces=true,
showspaces=false,
showstringspaces=false,
showtabs=false,
tabsize=2
}
\lstset{style=mystyle}

% autoref names
\def\makeautorefname#1#2{\csdef{#1autorefname}{#2}}
\makeautorefname{section}{Section}%
\makeautorefname{subsection}{Section}%
\makeautorefname{subsubsection}{Section}%

\begin{document}
\title[Additional material]{Additional material}
\author[D. W. Collison and D. Tubbenhauer]{Daniel W. Collison and Daniel Tubbenhauer}

\address{D.W.C.: The University of Sydney, School of Mathematics and Statistics F07, NSW 2006, Australia}
\email{daniel.collison@sydney.edu.au}

\address{D.T.: The University of Sydney, School of Mathematics and Statistics F07, Office Carslaw 827, NSW 2006, Australia, \href{http://www.dtubbenhauer.com}{www.dtubbenhauer.com}, https://orcid.org/0000-0001-7265-5047}
\email{daniel.tubbenhauer@sydney.edu.au}


\maketitle



\begin{lemma}\label{colorbasislem}
(See \cite[Lemma 6A.5]{penrose-cat}). Given the color basis, we have the following:
\begin{enumerate}
\item $p_k \cdot p_m = \delta_{km}p_m$, and hence $E_{k, 0} \cdot E_{m,0} = \delta_{km}E_{m, 0}$ and $E_{k, \pm} \cdot E_{m,\pm} = \delta_{km}E_{m, \pm}$, while all other products between color basis elements are zero;
\item $\Delta_{n, 1}(E_{k,0})= -\frac{1}{3} f_{n, k} E_{k,0} \otimes E_{k,0}$ and $\Delta_{n,1}(E_{k,\pm})= \frac{2}{3}f_{n,k} E_{k,\pm} \otimes E_{k,\pm}$;
\item $m_{n,1}(E_{k,0})= u g_{n,k}E_{k,0}$ and $m_{n,1}(E_{k,\pm})= (u+w) g_{n, k}E_{k,\pm}$,
\end{enumerate}
where 
\begin{align*}
f_{n, k} = \begin{cases*}
1 &  $n$ even \\
\omega^k &  $n$ odd
\end{cases*},
&  &
g_{n, k} = \begin{cases*}
\omega^{kn/2}&  $n$ even \\
\omega^{k(n+1)/2} &  $n$ odd
\end{cases*}.
\end{align*}
In other words, $\partial_{n,1}$ is diagonal in the color basis.
\end{lemma}

\begin{proof}
$\textbf{(a)}$ The calcuation is straightfoward using the following result:
$$p_k \,x^{n-\ell} = \omega^{-k\ell} p_k$$ for all $\ell = 1, \ldots, n$.

For $\textbf{(b)}$, let us use the fact that $\xi=1$ to simplify the expression for $\Delta_{n, 
1}$:
\begin{align*}
\Delta_{n,1}(v) =&  -\frac{1}{3n}  \sum\limits_{\substack{i+j \equiv c_n \,\text{mod} \, n }} (v \cdot x^i) \otimes x^j \\
&+ \frac{1}{3n}\sum\limits_{\substack{i+j \equiv (c_n-1) \,\text{mod} \, n }} (v \cdot x^i y^2) \otimes x^j + (v \cdot x^iy) \otimes x^jy + (v \cdot x^i) \otimes x^jy^2.
\end{align*}
for $c_n = \begin{cases*} 0 & $n$ even \\ 
1 & $n$ odd \end{cases*}$.

We now show that $\Delta_{n,1}(E_{k,0})= -\frac{1}{3}f_{n,k}E_{k,0} \otimes E_{k,0}$. The other cases involving $E_{k, \pm}$ proceed similarly. 

Firstly, 
\begin{align*}
E_{k,0} \otimes E_{k,0} 
&= p_k(1-x^{n-1}y^2) \otimes p_k(1-x^{n-1}y^2) \\
&= p_k \otimes p_k - \omega^{-k} p_k \otimes p_k y^2 - \omega^{-k} p_k y^2\otimes p_k + \omega^{-2k} p_k y^2\otimes p_k y^2 \\
&= \frac{1}{n^2}\sum_{r, s=0}^{n-1} \omega^{-k(r+s)}x^r \otimes x^s -\frac{1}{n^2}\sum_{r, s=0}^{n-1} \omega^{-k(r+s+1)}x^r \otimes x^sy^2 \\
&- \frac{1}{n^2}\sum_{r, s=0}^{n-1} \omega^{-k(r+s+1)}x^ry^2 \otimes x^s +\frac{1}{n^2}\sum_{r, s=0}^{n-1} \omega^{-k(r+s+2)}x^ry^2 \otimes x^sy^2 \\
&= \frac{1}{n^2}\sum_{\substack{\ell, r, s = 0 \\ r+s \equiv \ell \, \text{mod} \, n}}^{n-1} \omega^{-k\ell}x^r \otimes x^s -\frac{1}{n^2}\sum_{\substack{\ell, r, s = 0 \\ r+s +1\equiv \ell \, \text{mod} \, n}}^{n-1} \omega^{-k\ell}(x^r \otimes x^sy^2 + x^ry^2 \otimes x^s)\\
&+\frac{1}{n^2}\sum_{\substack{\ell, r, s = 0 \\ r+s +2\equiv \ell \, \text{mod} \, n}}^{n-1} \omega^{-k\ell}x^ry^2 \otimes x^sy^2.
\end{align*}
On the other hand,
\begin{align*}
\Delta_{n,1}(E_{k,0})&= \Delta_{n, 1}(p_k)-\omega^{-k}\Delta_{n,1}(p_ky^2
) \\
&=-\frac{1}{3n^2}\sum_{\substack{\ell, i, j = 0 \\ i+j\equiv c_n \, \text{mod} \, n}}^{n-1} \omega^{-k\ell}x^{\ell+i} \otimes x^j  +\frac{1}{3n^2}\sum_{\substack{\ell, i, j = 0 \\ i+j \equiv (c_n-1) \, \text{mod} \, n}}^{n-1} \omega^{-k\ell}x^{\ell+i}y^2 \otimes x^j  \\
&+ \frac{1}{3n^2}\sum_{\substack{\ell, i, j = 0 \\ i+j \equiv (c_n-1) \, \text{mod} \, n}}^{n-1} \omega^{-k\ell}( x^{\ell+i}y \otimes x^jy + x^{\ell+i} \otimes x^jy^2) \\&+\frac{1}{3n^2}\sum_{\substack{\ell, i, j = 0 \\ i+j\equiv c_n \, \text{mod} \, n}}^{n-1} \omega^{-k(\ell+1)}x^{\ell+i}y^2\otimes x^j - \frac{1}{3n^2}\sum_{\substack{\ell, i, j = 0 \\ i+j\equiv (c_n-1) \, \text{mod} \, n}}^{n-1} \omega^{-k(\ell+1)}x^{\ell+i+1}y^2\otimes x^j \\&- \frac{1}{3n^2}\sum_{\substack{\ell, i, j = 0 \\ i+j \equiv (c_n-1) \, \text{mod} \, n}}^{n-1} \omega^{-k(\ell+1)}( x^{\ell+i+1}y \otimes x^iy + x^{\ell+i}y^2 \otimes x^jy^2) 
\\
&=-\frac{1}{3n^2}f_{n,k}\sum_{\substack{\ell, r, s = 0 \\ r+s\equiv \ell \, \text{mod} \, n}}^{n-1} \omega^{-k\ell}x^{r} \otimes x^s  +\frac{1}{3n^2}f_{n,k}\sum_{\substack{\ell, r, s = 0 \\ r+s +1 \equiv \ell \, \text{mod} \, n}}^{n-1} \omega^{-k\ell}x^{r}y^2 \otimes x^s  \\
&+ \frac{1}{3n^2}f_{n,k}\sum_{\substack{\ell, r, s = 0 \\ r+s+1 \equiv \ell \, \text{mod} \, n}}^{n-1} \omega^{-k\ell} x^{r} \otimes x^sy^2 - \frac{1}{3n^2}f_{n,k}\sum_{\substack{\ell, r, s = 0 \\ r+s+2 \equiv \ell \, \text{mod} \, n}}^{n-1} \omega^{-k\ell}x^{r}y^2 \otimes x^sy^2,
\end{align*}
and so the result follows.

Finally, $\textbf{(c)}$ follows straightforwardly. For example, 
\begin{align*}
m_{n, 1}(E_{k,0})&=m_{n, 1}(p_k) - \omega^{-k}m_{n,1}(p_ky^2) \\
&=u p_k x^{n/2} + w p_k x^{(n-2)/2}y^2 - \omega^{-k}u p_k x^{n/2}y^2 - \omega^{-k}w p_k x^{(n-2)/2+1}y^2 \\
&= u p_k x^{n/2}- \omega^{-k}u p_k x^{n/2}y^2 
\\
& = u\omega^{kn/2}(p_k - \omega^{-k}p_ky^2) = u\omega^{kn/2} E_{k,0},
\end{align*}
for $n$ even, and the rest are similar.
\end{proof}

\begin{Theorem}\label{k32graphtheorem}
(See \cite[Theorem 7B.9]{penrose-cat}). For the complete graph $K_{4}$,
\begin{align*}
Mh_{n,1}(K_{4}) &= n\big(n^3(40b+41)-6n^2(13b+14) +11n(4b+5)-6(b+2)\big)t^4 \\& \hspace{1.5cm} + n\big(n^2(13b+14)-3n(4b+5) +2(b+2)\big)t^2,
\end{align*}
and for the complete bipartite graph $K_{3,2}$,
\begin{align*}
Mh_{n,1}(K_{3,2})= n\big(n^2(13b+14)-3n(4b+5) +2(b+2)\big)(t^6 +3t^4+3t^2+1).
\end{align*}
\end{Theorem}
\begin{proof}
We begin by inspecting the circle counts associated with the hypercube of states of the blowup of $K_{4}$, obtained using the code described in \cite{penrose-cat, CoTu-code}:
\begin{lstlisting}[language=Python]
circle counts by homological degree: [[2], [1, 2, 1, 1, 2, 1], [1, 2, 2, 1, 2, 1, 1, 3, 1, 2, 1, 2, 1, 2, 1], [2, 1, 2, 1, 3, 1, 3, 2, 3, 1, 1, 2, 1, 2, 2, 2, 1, 3, 2, 1], [2, 1, 2, 1, 2, 2, 2, 4, 2, 2, 2, 2, 1, 1, 2], [1, 3, 1, 1, 3, 1], [2]]
\end{lstlisting}
The maximum number of circles of any state is $4$, and so by using \cite[Remark 7B.2]{penrose-cat}, we only need to solve for the coefficients of a degree four polynomial in $n$ without a constant term,
\begin{align*}
c_1n^4 + c_2n^3 + c_3n^2 + c_4n,
\end{align*}
for each degree of $t$ up to six and $b$ up to one, using the data from \autoref{table:completepoly} for $K_{4}$ up to $n=4$. The only nontrivial systems of equations that require solving are 
\begin{align*}
\alpha_1 + \alpha_2 + \alpha_3 + \alpha_4 &=0 & \beta_1 + \beta_2 + \beta_3 + \beta_4 &=0 \\
16 \alpha_1 + 8 \alpha_2 + 4\alpha_3 + 2\alpha_4 &= 180 & 16 \beta_1 + 8 \beta_2 + 4\beta_3 + 2\beta_4 &= 180 \\
81 \alpha_1 + 27 \alpha_2 + 9\alpha_3 + 3\alpha_4 &= 1512 & 81 \beta_1 + 27 \beta_2 + 9\beta_3 + 3\beta_4 &= 1512 \\
256 \alpha_1 + 64 \alpha_2 + 16\alpha_3 + 4\alpha_4 &= 5928 & 256 \beta_1 + 64 \beta_2 + 16\beta_3 + 4\beta_4 &= 5952
\end{align*}
associated with degree one and zero in $b$ respectively for degree four in $t$, and similarly for degree two in $t$:
\begin{align*}
 \delta_1 + \delta_2 + \delta_3 + \delta_4 &=3  & \gamma_1 + \gamma_2 + \gamma_3 + \gamma_4 &=3\\
 16 \delta_1 + 8 \delta_2 + 4\delta_3 + 2\delta_4 &= 60 & 
16 \gamma_1 + 8 \gamma_2 + 4\gamma_3 + 2\gamma_4 &= 60  \\
81 \delta_1 + 27 \delta_2 + 9\delta_3 + 3\delta_4 &= 249 & 
81 \gamma_1 + 27 \gamma_2 + 9\gamma_3 + 3\gamma_4 &= 255 \\
256 \delta_1 + 64 \delta_2 + 16\delta_3 + 4\delta_4 &= 648 & 
256 \gamma_1 + 64 \gamma_2 + 16\gamma_3 + 4\gamma_4 &= 672 
\end{align*}
The reader is encouraged to check that the solutions are given by 
\begin{align}\label{solutions}
\begin{pmatrix}
\alpha_1 \\
\alpha_2 \\ 
\alpha_3 \\
\alpha_4
\end{pmatrix} &= \begin{pmatrix}
40 \\
-78 \\ 
44 \\
-6
\end{pmatrix}, & \begin{pmatrix}
\beta_1 \\
\beta_2 \\ 
\beta_3 \\
\beta_4
\end{pmatrix} &= \begin{pmatrix}
41 \\
-84 \\ 
55 \\
-12
\end{pmatrix}, & \begin{pmatrix}
\delta_1 \\
\delta_2 \\ 
\delta_3 \\
\delta_4
\end{pmatrix} &= \begin{pmatrix}
0 \\
13 \\ 
-12 \\
2
\end{pmatrix}, & \begin{pmatrix}
\gamma_1 \\
\gamma_2 \\ 
\gamma_3 \\
\gamma_4
\end{pmatrix} &= \begin{pmatrix}
0 \\
14 \\ 
-15 \\
4
\end{pmatrix}.
\end{align}

For $K_{3, 2}$, the circle counts associated with the hypercube of states of the blowup are given by:
\begin{lstlisting}[language=Python]
circle counts by homological degree: [[3], [2, 2, 2, 2, 2, 2], [3, 1, 1, 1, 1, 1, 1, 1, 1, 3, 1, 1, 1, 1, 3], [2, 2, 2, 2, 2, 1, 1, 1, 1, 2, 2, 1, 1, 1, 1, 2, 2, 2, 2, 2], [3, 1, 1, 1, 1, 3, 1, 1, 1, 1, 1, 1, 1, 1, 3], [2, 2, 2, 2, 2, 2], [3]]
\end{lstlisting}
As above, we only need to solve for the coefficients of a degree three polynomial in $n$ without a constant term, \ie\, of the form $
c_1n^3 + c_2n^2 + c_3n$, for each degree of $t$ up to six and $b$ up to one, using the data from \autoref{table:bicompletepoly} for $K_{3, 2}$ up to $n=3$. Comparing with the data for $K_4$, it is straightforward to see that in fact the solutions for the coefficients $\delta_i$, $\gamma_i$, $i =1, \ldots, 4$ from \autoref{solutions} can be reused to solve the nontrivial systems of equations for $K_{3, 2}$, up to overall constant factors.
\end{proof}


\begin{table}
\renewcommand{\arraystretch}{1}
\centering
\begin{tabular}{| c || w{c}{.75cm} | c| c| }
\hline 
$4$ &  $0$ & $4(15 b + 18)(3t^2 + 1)$ & $24(247b+248)t^4 + 24(27b + 28)t^2$   \\ \hline
$3$ &  $0$  & $3(11 b + 13)(3t^2 + 1)$ & $504(3b+3)t^4 + 3(83b + 85)t^2$ \\ \hline
$2$ &  $0$  & $2(7b + 8)(3t^2 + 1)$ & $20(3b + 3)(3t^4+t^2)$ \\ \hline
$1$ &  $0$ & $(3 b + 3)(3t^2 + 1)$ & $(3b+3)t^2$\\\hline \hline
\diagbox[dir=SW]{$n$}{$m$} & $2$& $3$ & $4$ \\ \hline
\end{tabular} 
\caption{$Mh_{n,1}(K_m)$ for a range of $m, n$.}
\label{table:completepoly}
\end{table}
\begin{table}
\renewcommand{\arraystretch}{1}
\centering
\begin{tabular}{| c || w{c}{.75cm} | w{c}{.75cm} | c| c| }
\hline 
$4$ &  $0$ & $0$ & $4(15 b + 18)(t^4 + 6t^2 + 1)$ & $24(27b+28)(t^6 + 3t^4+ 3t^2 +1)$  \\ \hline
$3$ &  $0$ & $0$ & $3(11 b + 13)(t^4 + 6t^2 + 1)$ & $3(83b+85)(t^6 + 3t^4+ 3t^2 +1)$ \\ \hline
$2$ &  $0$ & $0$ & $2(7 b + 8)(t^4 + 6t^2 + 1)$ & $20(3b+3)(t^6 + 3t^4+ 3t^2 +1)$  \\ \hline
$1$ &  $0$ & $0$ & $(3b + 3)(t^4+6t^2 + 1)$ & $(3b+3)(t^6 + 3t^4+ 3t^2 +1)$ \\\hline \hline
\diagbox[dir=SW]{$n$}{$(m, p)$} & $(1, 1)$& $(2,1) $& $(2,2)$ & $(3, 2)$ \\ \hline
\end{tabular} 
\caption{$Mh_{n,1}(K_{m,p})$ for a range of $m, p,  n$.}
\label{table:bicompletepoly}
\end{table}




%%%%%%%%%%%%%%%%%%%%%%%%%%%%%%%%%%%%%%%%%%%%%
%\bibliographystyle{alphaurl}
%\bibliography{references}
%%%%%%%%%%%%%%%%%%%%%%%%%%%%%%%%%%%%%%%%%%%%%


\begin{thebibliography}{BHMV95}


\bibitem[CT26a]{penrose-cat}
D.~W. Collison and D.~Tubbenhauer.
\newblock Categorifying the Penrose polynomial using nonorientable surfaces.
\newblock 2026.
\newblock \url{https://arxiv.org/abs/2607.01632}.


\bibitem[CT26b]{CoTu-code}
D.~W. Collison and D.~Tubbenhauer.
\newblock Code, data and more for ``Categorifying the Penrose polynomial using nonorientable surfaces''.
\newblock 2026.
\newblock \url{https://github.com/dtubbenhauer/graphhomology}.


\end{thebibliography}

\end{document}